% Every result number is bound; j.x2a.* and j.x5.* come from
% values/journal-transcribed.json, and j.x3.* from the tabulations of the
% released free-share catalog (rounds 4 and 5).
Opinion-dynamics models explain collective patterns through rules for how individuals revise their views in contact with others \cite{flache2017}; we use attitude and opinion interchangeably for those views. Recent work calibrates such models to individual transitions in panel surveys: Gestefeld and Lorenz fit a model with interpersonal updating to eleven annual waves of trust in government in the Swiss Household Panel (SHP) \cite{gestefeld2023}. Panels, however, record what people answer, not whom they talk to; the SHP observes no exposure beyond household co-residence. A model that reproduces short-term stability may fit for reasons unrelated to influence, so a calibration needs further evidence before its fit can count in favor of the interaction mechanism.

This article investigates the limits of discriminating network influence from individual persistence in attitude panels. The central issue is what a difference between their observed trajectories can establish when contacts are unobserved, measurement is imperfect and persistence can reproduce the same short-term stability. We quantify one aspect of that problem with a detection boundary: the nominal influence input at which a fixed classifier separates two specified panel distributions, as a function of response reliability. The boundary summarizes a conditional comparison; its meaning depends on the measurement assumptions and the persistence alternatives allowed.

The question is hard for two reasons. The first is measurement: a changed answer can reflect real change, a passing state or reporting error. Classical panel models separate the stability of a construct from the reliability of its measurement \cite{heise1969}, vague questions can make attitudes look more volatile than they are \cite{achen1975}, and updating between every wave may be hard to tell from settled dispositions measured with error \cite{kiley2020}. The second is identification. Observed behavior cannot reveal whether a group's average behavior affects its members unless reference groups are known \cite{manski1993}, and in observational network data contagion and homophily are generically confounded \cite{shalizi2011}. Without exposure data, influence from unobserved contacts and individual persistence can produce the same observed transitions; where they do not, the difference must appear somewhere in the trajectories.

Two analytic observations (Section~\ref{sec:methods-identification}) make the distinction concrete. Linear assimilation toward the population mean is observationally equivalent to individual autoregression matched at lag one: the population area under the receiver operating characteristic curve (AUC), the probability that a panel with influence is ranked above one without, is 0.5. For fixed random contacts, the influence term instead contains a contribution from the person's own state and an input from persistent contacts. The emphasis on the serial correlation of that input, and its possible role in discrimination, is an interpretation developed after the results became available.

% Design schematic (Fig. 1). Needs \usepackage{tikz} and
% \usetikzlibrary{arrows.meta,positioning,calc} in the preamble.
% The inset curve is a drawing, not data. Text sizes 7pt and 8pt bold use
% fonts already in the offline Tectonic cache.
\begin{figure*}[t]
\centering
\begin{tikzpicture}[
  font=\fontsize{7pt}{8.3pt}\selectfont,
  >={Stealth[length=1.5mm,width=1.3mm]},
  box/.style={draw=darkgray, rounded corners=1.5pt, align=flush left,
    inner sep=3pt, text width=#1},
  box/.default=2.2cm,
  flow/.style={->, semithick, darkgray},
]
\node[box=3.45cm, fill=blue!7] (influence)
  {{\fontsize{8pt}{9.4pt}\selectfont\bfseries Influence arm}\\[1.5pt]
   Latent attitudes persist between waves. Each person is pulled toward
   contacts within distance~1 and pushed away from farther ones. Eight
   fixed random contacts per person, never observed. Nominal
   input~$a$: the variance fraction assigned to contacts in the innovation.};
\node[box=3.45cm, fill=orange!12, below=7.5mm of influence] (persistence)
  {{\fontsize{8pt}{9.4pt}\selectfont\bfseries Persistence model, no network}\\[1.5pt]
   A~autoregressive; B~trait-plus-autoregressive (main comparison);
   C~autoregressive with an altered response channel; D~state-dependent;
   E~mover--stayer.};
\draw[->, densely dashed, darkgray]
  ($(influence.south west)+(3.5mm,0)$) --
  node[right, align=flush left, inner sep=2pt] {matched to its short-lag\\
    correlations}
  ($(persistence.north west)+(3.5mm,0)$);

\coordinate (between) at ($(influence.east)!0.5!(persistence.east)$);
\node[box=2.35cm, fill=black!5, right=7mm of between] (channel)
  {{\fontsize{8pt}{9.4pt}\selectfont\bfseries Response channel}\\[1.5pt]
   Same for both arms, except in comparisons with model~C. Signal
   share~$r$, known. Answers in eleven ordinal categories, assigned within
   each wave. Same attrition.};
\coordinate (join) at ($(channel.west)+(-3mm,0)$);
\draw[semithick, darkgray] (influence.east) -| (join);
\draw[semithick, darkgray] (persistence.east) -| (join);
\draw[flow] (join) -- (channel.west);

\node[box, right=4mm of channel] (features)
  {{\fontsize{8pt}{9.4pt}\selectfont\bfseries Panel features}\\[1.5pt]
   Correlations of each person's answers 3, 4 and 5 waves apart, pooled
   over people and waves.};
\node[box, right=4mm of features] (classifier)
  {{\fontsize{8pt}{9.4pt}\selectfont\bfseries Fixed classifier}\\[1.5pt]
   Linear discriminant trained on 2,000 panels per arm. It scores whole
   panels.};
\draw[flow] (channel) -- (features);
\draw[flow] (features) -- (classifier);

\node[box=2.75cm, right=4mm of classifier, minimum height=3.55cm]
  (boundary) {};
\node[anchor=north west, inner sep=3pt, text width=2.75cm,
  align=flush left] at (boundary.north west)
  {{\fontsize{8pt}{9.4pt}\selectfont\bfseries Boundary $a^*(r)$}\\[1.5pt]
   First of the nine grid strengths at which the expected AUC
   reaches 0.70.};
\draw[flow] (classifier) -- (boundary);
% Inset drawn in centimeters; the AUC axis runs from 0.45 at the origin.
\begin{scope}[shift={($(boundary.south west)+(0.62cm,0.42cm)$)}]
  \fill[blue!12] (0.000cm,0.045cm) -- (0.240cm,0.068cm) -- (0.480cm,0.090cm) -- (0.720cm,0.113cm) -- (0.960cm,0.202cm) -- (1.200cm,0.585cm) -- (1.440cm,0.922cm) -- (1.680cm,1.103cm) -- (1.920cm,1.147cm) -- (1.920cm,1.238cm) -- (1.680cm,1.238cm) -- (1.440cm,1.147cm) -- (1.200cm,0.833cm) -- (0.960cm,0.360cm) -- (0.720cm,0.225cm) -- (0.480cm,0.202cm) -- (0.240cm,0.180cm) -- (0.000cm,0.180cm) -- cycle;
  \draw[->, darkgray] (0cm,0cm) -- (2.088cm,0.000cm) node[right, inner sep=1pt] {$a$};
  \draw[->, darkgray] (0cm,0cm) -- (0.000cm,1.350cm);
  \node[rotate=90, anchor=south, inner sep=1pt] at (-0.08cm,1.0cm) {AUC};
  \draw[dashed, gray] (0.000cm,0.562cm) -- (2.040cm,0.562cm);
  \node[left, inner sep=1.5pt] at (0.000cm,0.562cm) {0.70};
  \draw[darkgray] (0.000cm,0.113cm) -- (0.240cm,0.124cm) -- (0.480cm,0.146cm) -- (0.720cm,0.169cm) -- (0.960cm,0.281cm) -- (1.200cm,0.709cm) -- (1.440cm,1.035cm) -- (1.680cm,1.170cm) -- (1.920cm,1.193cm);
  \fill[darkgray] (0.000cm,0.113cm) circle (0.75pt);
  \fill[darkgray] (0.240cm,0.124cm) circle (0.75pt);
  \fill[darkgray] (0.480cm,0.146cm) circle (0.75pt);
  \fill[darkgray] (0.720cm,0.169cm) circle (0.75pt);
  \fill[darkgray] (0.960cm,0.281cm) circle (0.75pt);
  \fill[darkgray] (1.200cm,0.709cm) circle (0.75pt);
  \fill[darkgray] (1.440cm,1.035cm) circle (0.75pt);
  \fill[darkgray] (1.680cm,1.170cm) circle (0.75pt);
  \fill[darkgray] (1.920cm,1.193cm) circle (0.75pt);
  \draw[densely dotted, darkgray] (1.200cm,0.709cm) -- (1.200cm,0.000cm)
    node[below, inner sep=1.5pt] {$a^*(r)$};
\end{scope}
\end{tikzpicture}
\caption{Design of the simulation. At each signal share $r$ and influence
strength $a$, panels from the influence arm and from a persistence model
matched to its short-lag correlations pass through the same response channel
(model C alters it in the persistence arm). A fixed classifier scores whole
panels from their long-lag correlations. The boundary $a^*(r)$ is the first
of the nine grid strengths at which the expected AUC reaches 0.70; the shaded area stands for the simultaneous uncertainty
band. Matching, training and testing use separate simulated panels. The inset
is a drawing, not data. \esmresource, Table~\ref{esm-tab:glossary-terms}, defines
the terms used in the article.}
\label{fig:design}
\end{figure*}


We study these limits in a prespecified simulation in which the generating mechanism is known (Fig.~\ref{fig:design}). Synthetic eleven-wave panels approximate the size and refreshment design of the SHP trust series. In the influence arm, persistent latent attitudes are pulled toward similar contacts and pushed away from dissimilar ones on an unobserved random network. The parameter $a$ sets the nominal fraction of innovation variance assigned to contacts; it is neither the realized fraction of attitude change attributable to them nor a directly interpretable measure of social influence in a survey. Five persistence alternatives are matched to short-lag population correlations. The response channel has known signal share $r$, its reliability before discretization. A fixed linear classifier scores whole panels from their pooled correlations three to five waves apart, and $a^*(r)$ is the first grid strength at which the expected AUC reaches \resultnum{v4.protocol.target}{0.70}.

The contribution is a quantitative account of the conditions under which such discrimination is informative and where it breaks down. The design separates successful comparisons from failures to fit or validate, rather than treating a failed alternative as evidence for influence. It also distinguishes the prespecified results from later exploratory analyses that examine how the result changes when measurement and persistence assumptions are relaxed. The protocol was adopted after partial exposure to calibration results, and the confirmatory run used new seeds; the study is prespecified rather than publicly registered (Section~\ref{sec:methods-record}).

Against trait-plus-autoregressive persistence, the estimated AUC first reached \resultnum{v4.protocol.target}{0.70} at $a=\resultnum{j.b.r055.first.point}{0.10}$ for every signal share from \resultnum{v4.tab.T1.B.r055.r}{0.55} to \resultnum{v4.tab.T1.B.r090.r}{0.90} and at $a=\resultnum{j.b.r050.first.point}{0.30}$ for $r=\resultnum{v4.tab.T1.B.r050.r}{0.50}$. Under an equal-variance normal approximation at a one-sided 5\% level, an AUC of 0.70 gives a test on a single panel a power of about \resultnum{j.auc_translation.target070.power}{0.18} against a known persistence score distribution.

These confirmatory findings establish discrimination against the main persistence alternative under favorable measurement assumptions. The prespecified comparison across all persistence models could not be interpreted, because the altered response channel failed the zero-strength validity rule and the declared solver did not obtain acceptable state-dependent fits near the main boundary. The exploratory analyses then show why the main result is sensitive to the comparison design. A changing response channel separated panels without influence, while allowing persistence to fit reliability and an additional lag absorbed most separation by a classifier using the remaining lags. The fitted reliability differed from the design value by only \resultnum{j.x3.free.r_shift.min}{0.0011}--\resultnum{j.x3.free.r_shift.max}{0.0028}; this combined change does not isolate the effect of reliability alone.

The resulting distinction between discrimination and attribution is consequential for models calibrated to attitude panels \cite{manski1993,shalizi2011,cohencole2008}. The Discussion examines three limits: the restricted persistence alternatives, privileged knowledge of reliability and the absence of network burn-in. It separates sensitivities demonstrated by the exploratory analyses from the untested effects of richer persistence and initial transients, and considers what additional evidence would make an influence interpretation more credible (Section~\ref{sec:implications}).
